\chapter{Innledning}
\label{ch:intro}
\section{Problem og mål}
Oppgaven i VE3140 ber om en fungerende sanntidsprototype i Ada på micro:bit V2
og Mecanum-bil, en LaTeX-rapport med teori og deadlineanalyse, vesentlig
Ada-kode som vedlegg og en vurdering av prosjektgjennomføringen. Den bevarte
oppgaveteksten i \kode{grunnlag/oppgavetekst.png} anbefaler 15--20 sider
rapporttekst og oppgir at kodevedlegg kommer i tillegg. Rapporten følger
den opprinnelige strukturen med innledning, metode, resultater og diskusjon.

Problemstillingen er: \emph{Hvordan kan micro:bit V2 kombinere avstandsmåling,
kjørebeslutninger og motorstyring slik at bilen beveger seg autonomt og
reagerer forutsigbart på hindringer og feil?} En tilfredsstillende løsning
må både virke fysisk og gjøre reaksjonstiden etterprøvbar. Kompilering og
et høyt antall tester alene besvarer ikke dette.

Målet for demonstrasjonen er fremoverkjøring, stopp ved hindring,
stasjonær retningsskanning, klarert sving og ny fremoverkjøring. Det skal
være mulig å forklare hvilken tilstand bilen befinner seg i, hvorfor en
måling godtas eller avvises, og hva som skjer når en komponent svikter.
Bremselys og blinklys skal formidle styringens intensjon.

\section{Avgrensning og dokumentgrunnlag}
Arbeidet gjelder en liten undervisningsrobot i et kontrollert område.
Raspberry Pi er utviklingsvert; autonomien avhenger ikke av en kontinuerlig
USB-forbindelse. Koden bruker faste PWM-kommandoer, uten målt hjulhastighet
eller lukket regulering. Kartbygging, dynamiske hindringer, PID og
servoskanning mens bilen kjører er utenfor den validerte løsningen.

Gjeldende kode, rålogger, den bevarte oppgaven og produsentdokumentasjon er
kontrollgrunnlaget. Forelesningsfilen \kode{ISforelesninger.pdf} og en
gjeldende sensurveiledning var ikke tilgjengelige ved denne gjennomgangen.
Tidligere oppgitte forelesningssider brukes derfor ikke som ny kontrollert
evidens. Faglig sporbarhet presenteres i kravtabellene; full pensumdekning
må avstemmes når disse kildene er tilgjengelige.

\section{KI-bruk og forfatteransvar}
I denne oppfølgingen er Codex brukt til kodegjennomgang, kodeendringer,
testutforming, bevisoppsett og redigering av LaTeX og presentasjon.
Programtester og bygg er kjørt på Pi-en. Fysisk oppsett og observasjoner
kommer fra brukerens opplysninger og daterte logger.
\todo{Gruppen må beskrive øvrig KI-bruk, kontrollere denne deklarasjonen
mot emnets regler og selv kunne forklare og forsvare koden.}
Navn og studentnumre er videreført fra originalrapporten; arbeidsfordeling
og ikke-dokumenterte gruppemøter er ikke konstruert.
